% ---
% id: protocol-compiler/rem:summation-condition
% kind: rem
% title: "The summation condition beside completeness"
% summary: >-
%   Completeness alone says that the computed floating-point sum is the
%   identity; the summation condition is what makes it an identity of the
%   declared numbers. A family of exactly representable entries may therefore
%   still be refused, since the condition is on the additions.
% prominence: standard
% status: accepted
% uses: [def:exact]
% mentions: [foundations/rem:computed-completeness, foundations/lem:exact-accumulation]
% assumes: []
% refs: []
% provenance: newly-derived
% lean:
% ---

\begin{remark}\label{rem:summation-condition}
The total $\Sigma^{(s)}$ of Definition~\ref{def:exact} is by
\eqref{eq:declorder} a \emph{floating-point} sum. So condition (C) alone says
that the \emph{computed} sum is the identity, and \emph{not} that
$\sum_ja_j=\idon{\Qsys}$ holds of the declared numbers, which is what every
derivation of this contract that consumes completeness reads it as, and what
a consumer of the emitted declaration reads it as. Condition (S) is exactly
what makes the two the same statement
(\xref{foundations/lem:exact-accumulation}), and it is the reason a family
whose entries are all exactly representable may still be refused: the
condition is on the \emph{additions}, not on the operands.
\xref{foundations/rem:computed-completeness} exhibits a family satisfying
(H), (P) and completeness of the computed sum whose true sum is not the
identity.
\end{remark}
